% Current v15 support consequences. All graph counts are necessary candidates.
\subsection{三角形を必要とする高分母 support と残る十三型}
\label{sf15:section}

\begin{corollary}[三角形の必要性]
\label{sf15:triangle-required}
正の局所整数 charts，marked puncture の整数 clutching，対向 $D$ の相殺，
および node 非負性を満たす，五次の標準 $K_5$ branch graph を考える．
高分母 support $H=\{ij:q_{ij}>1\}$ が非空なら，$H$ は三角形を含む．
従って $H$ が triangle-free なら，全ての分母は一である．
\end{corollary}

\begin{proof}
系~\ref{fo13:remaining-support} により，辺を持つ成分は一つで，
四または五頂点を使い，閉路を含む．$H$ に三角形がないと仮定する．
五角形を含む場合，その対角線はどれも三角形を作るので，
$H=C_5$ となり，系~\ref{cf13:degree-two-rigidity} に反する．

従って四角形を含む．その二本の対角線も三角形を作るため加えられない．
第五頂点から四角形への高辺が二本以上あれば，その行先は隣接できないので，
高辺はちょうど二本で，行先は対向する二頂点である．
従って可能な形は，$C_4$，$C_4$ に葉辺を一本付けた形，$K_{2,3}$ の三つに尽くされる．
これらはそれぞれ系~\ref{cf13:degree-two-rigidity}，
定理~\ref{cp14:no-pendant-square}，定理~\ref{kb15:no-K23} に反する．
従って $H$ は三角形を含む．
\end{proof}

\begin{corollary}[残る support の構造と個数]
\label{sf15:remaining-support}
前系の仮定の下で，高分母 support $H$ が非空なら，辺を持つ連結成分は一つで，
四または五頂点を使い，少なくとも五辺を持つ．最大次数は三以上である．
四頂点を使う場合には $H\cong K_4$ に限られる．
唯一の閉路を持つ場合には，三角形の異なる二頂点に葉辺を一本ずつ付けた graph に限られる．

ここまでの graph 障害を通過する型は，標識を忘れると十三，
$K_5$ の標識を保つと四百二十一である．
これは必要条件を通過した graph の個数であり，整数 charts の接着や幾何的実現の存在を主張しない．
\end{corollary}

\begin{proof}
連結成分と最大次数の主張は系~\ref{fo13:remaining-support} による．
四頂点を使う connected graph で四辺以下なら，木，$C_4$，または三角形に葉辺を付けた形である．
木と $C_4$ は既に排除され，最後の形は三角形に二つの二価頂点を持つので，
定理~\ref{tn14:no-degree-two-triangle} に反する．
五辺の場合は diamond であり，定理~\ref{dg15:no-diamond} で排除された．
残る六辺の場合が $K_4$ である．五頂点を使う場合には閉路を含むので，辺数は五以上である．

唯一の閉路を持つ場合，前系によってそれは三角形である．
三角形に長さ二の道を付ける場合，または同じ頂点に二本の葉辺を付ける場合には，
三角形の二頂点が二価のままなので，定理~\ref{tn14:no-degree-two-triangle} に反する．
従って三角形の異なる二頂点に葉辺を一本ずつ付けた形だけが残る．

最後の個数は，$K_5$ の全 $2^{10}=1024$ 部分 graph に対する有限な検算で得られる．
木成分・単純閉路成分，二価二頂点を持つ三角形，葉辺付き四角形を除いた十五型から，
diamond と $K_{2,3}$ を除く．両者の標識付き個数はそれぞれ三十と十なので，
型数は十五から十三，標識付き個数は四百六十一から四百二十一となる．
全 graph の列挙，頂点置換による同型判定，および各代表の自己同型群の検算は
\nolinkurl{verify_support_agenda_v15.py} と
\nolinkurl{support_agenda_verification_v15.json} に保存した．
この検算は分母や行列 entry の有界探索ではない．
\end{proof}

二価 core に限れば，系~\ref{kb15:negative-closes-triangle} がさらに局所的な制約を与える．
その core の二つの高分母近傍が互いに低辺で結ばれるなら，二高分母標準形は正の枝に属する．
この局所制約自体は対向 $D$ の相殺を使わない．

唯一残る unicyclic 型では，定理~\ref{tl15:remaining-positive} により，
中央の二価 core は交差した低 ray の対応を持つ正の枝であり，
$r\leq4$ または $z=1$ に限られる．この帰着は，その support 全体の排除を意味しない．
